\documentclass[11pt]{article}
\usepackage{hw-solutions-common}

\begin{document}

\hwsolhead{Written HW 6}{Problem statements from Knight, \emph{Physics for Scientists and
Engineers} (5th ed., Global Edition), Chapter 28. $g = 9.8$ m/s$^2$ throughout.}

%============================================================
\problem{28.17}{3}{The voltage across the terminals of a 9.0 V battery is 8.5 V when the
battery is connected to a 20 $\Omega$ load. What is the battery's internal resistance?}

\model{A real battery is an ideal emf $\varepsilon$ in series with a small internal
resistance $r$. The terminal voltage is whatever's left of $\varepsilon$ after the
current drops some of it across $r$.}

\visualize{The load current is set by the terminal voltage across the external
resistor, $I=\Delta V_{\text{bat}}/R$. The ``missing'' $\varepsilon-\Delta
V_{\text{bat}}=0.5$ V is what gets dropped across $r$, by that same current.}

\solve{
\[
I = \frac{\Delta V_{\text{bat}}}{R} = \frac{8.5\text{ V}}{20\ \Omega} = 0.425\text{ A}
\]
\[
r = \frac{\varepsilon-\Delta V_{\text{bat}}}{I} = \frac{9.0\text{ V}-8.5\text{ V}}{0.425\text{
A}} = \boxed{1.2\ \Omega}.
\]}

\review{An internal resistance around 1 $\Omega$ is typical for a battery this size, and
it's exactly enough to knock the terminal voltage down by the stated 0.5 V at this load
current --- consistent with the numbers given.}

%============================================================
\problem{28.20}{3}{A variable resistor $R$ is connected across the terminals of a
battery. Figure EX28.20 shows the current in the circuit as $R$ is varied: the graph
passes through roughly $(0\ \Omega, 6.0\text{ A})$ and $(10\ \Omega, 3.0\text{ A})$.
What are the emf and internal resistance of the battery?}

\model{Same real-battery loop rule as above, now solved the other direction: given
$I$ at two different values of $R$, find $\varepsilon$ and $r$.}

\visualize{The loop rule gives $\varepsilon = I(R+r)$ for any point on the curve. Reading
off two points turns this into two equations in the two unknowns $\varepsilon$ and $r$.}

\solve{At $R=0\ \Omega$: $\varepsilon = (6.0\text{ A})\,r$. At $R=10\ \Omega$: $\varepsilon
= (3.0\text{ A})(r+10\ \Omega)$. Setting these equal,
\[
(6.0\text{ A})\,r = (3.0\text{ A})(r+10\ \Omega)
\;\Rightarrow\;
2r = r+10\ \Omega
\;\Rightarrow\;
\boxed{r = 10\ \Omega},
\]
\[
\varepsilon = (6.0\text{ A})(10\ \Omega) = \boxed{60\text{ V}}.
\]}

\review{Check: at $R=10\ \Omega$, $I=\varepsilon/(R+r)=60\text{ V}/20\ \Omega=3.0$ A, which
matches the second point exactly, and the curve $I(R)=60\text{ V}/(R+10\ \Omega)$
decreases smoothly and concavely as $R$ grows, just like the figure.}

%============================================================
\problem{28.33}{2}{What is the time constant for the discharge of the capacitors in
Figure EX28.33? The circuit has two 2 $\mu$F capacitors and two 1 k$\Omega$ resistors,
all four connected directly across the same two rails (all in parallel with each
other).}

\model{With everything on the same two rails, the whole network reduces to one
equivalent capacitor discharging through one equivalent resistor.}

\visualize{Capacitors in parallel add directly; resistors in parallel combine by
reciprocals.}

\solve{
\[
C_{\text{eq}} = 2\ \mu\text{F}+2\ \mu\text{F} = 4\ \mu\text{F}, \qquad
R_{\text{eq}} = \left(\frac{1}{1\text{ k}\Omega}+\frac{1}{1\text{
k}\Omega}\right)^{-1} = 500\ \Omega
\]
\[
\tau = R_{\text{eq}}C_{\text{eq}} = (500\ \Omega)(4\times10^{-6}\text{ F}) =
\boxed{2.0\text{ ms}}.
\]}

\review{A millisecond-scale time constant is typical for $\mu$F capacitors discharging
through kilohm-scale resistors --- the right order of magnitude.}

%============================================================
\problem{28.34}{2}{What is the time constant for the discharge of the capacitors in
Figure EX28.34? Now the two 2 $\mu$F capacitors are in series with each other, and the
two 1 k$\Omega$ resistors are in series with each other, around a single loop.}

\model{Same idea as 28.33, but now both the capacitors and the resistors combine
\emph{in series} instead of in parallel.}

\visualize{Capacitors in series combine by reciprocals; resistors in series add
directly --- the opposite pairing from the previous problem.}

\solve{
\[
\frac{1}{C_{\text{eq}}} = \frac{1}{2\ \mu\text{F}}+\frac{1}{2\ \mu\text{F}}
\;\Rightarrow\;
C_{\text{eq}} = 1\ \mu\text{F}, \qquad
R_{\text{eq}} = 1\text{ k}\Omega+1\text{ k}\Omega = 2\text{ k}\Omega
\]
\[
\tau = R_{\text{eq}}C_{\text{eq}} = (2000\ \Omega)(1\times10^{-6}\text{ F}) =
\boxed{2.0\text{ ms}}.
\]}

\review{This comes out to the \emph{same} 2.0 ms as the parallel arrangement in 28.33 ---
but that's a coincidence of these particular numbers (the resistor combination
quadruples while the capacitor combination quarters, and $4\times\frac14=1$), not a sign
that series and parallel circuits behave the same way in general.}

%============================================================
\problem{28.36}{4}{The switch in Figure EX28.36 has been in position 1 for a long time.
It is changed to position 2 at $t=0$ s. The battery is 9 V, the capacitor is 4
$\mu$F, and the resistor is 25 $\Omega$. What are the charge $Q$ on the capacitor and
the current $I$ through the resistor (a) immediately after the switch is closed? (b) at
$t=50\ \mu$s? (c) at $t=200\ \mu$s?}

\model{In position 1, the capacitor has had unlimited time to reach steady state, so no
current flows into it and it sits at the full battery voltage. Moving the switch to
position 2 disconnects the battery and leaves the capacitor discharging through the
resistor alone --- a standard $RC$ discharge, $Q(t)=Q_0e^{-t/\tau}$.}

\visualize{The capacitor's initial charge is set by the battery emf before the switch
flips: $Q_0=C\varepsilon$. After the switch flips, both $Q$ and $I=I_0e^{-t/\tau}$ decay
with the same time constant $\tau=RC$.}

\solve{(a) Immediately before (and so also immediately after, since charge can't jump
instantaneously) the switch moves,
\[
Q_0 = C\varepsilon = (4\times10^{-6}\text{ F})(9\text{ V}) = \boxed{36\ \mu\text{C}},
\qquad
I_0 = \frac{\Delta V_{\text{R}}}{R} = \frac{9\text{ V}}{25\ \Omega} = \boxed{0.36\text{
A}}.
\]

(b)--(c) The time constant is
\[
\tau = RC = (25\ \Omega)(4\times10^{-6}\text{ F}) = 100\ \mu\text{s}.
\]
At $t=50\ \mu$s ($t/\tau=0.50$):
\[
Q = (36\ \mu\text{C})e^{-0.50} = \boxed{22\ \mu\text{C}}, \qquad
I = (0.36\text{ A})e^{-0.50} = \boxed{0.22\text{ A}}.
\]
At $t=200\ \mu$s ($t/\tau=2.0$):
\[
Q = (36\ \mu\text{C})e^{-2.0} = \boxed{4.9\ \mu\text{C}}, \qquad
I = (0.36\text{ A})e^{-2.0} = \boxed{49\text{ mA}}.
\]}

\review{$Q$ and $I$ decay together, both by the same factor $e^{-t/\tau}$ at every
instant, since $I=dQ/dt$ is just proportional to $Q$ for a simple $RC$ discharge --- by
$t=200\ \mu$s (two full time constants), both are down to about 14\% of their starting
values, exactly $e^{-2.0}$.}

%============================================================
\problem{28.49}{4}{Suppose you have resistors 2.5 k$\Omega$, 3.5 k$\Omega$, and 4.5
k$\Omega$ and a 100 V power supply. What is the ratio of the total power delivered to
the resistors if they are connected in parallel to the total power delivered if they are
connected in series?}

\model{In both arrangements the \emph{same} 100 V is applied across the whole
combination, so the total power delivered is $P=(\Delta V)^2/R_{\text{eq}}$ in each case
--- the ratio of powers is just the inverse ratio of the two equivalent resistances.}

\visualize{A parallel combination's equivalent resistance is always \emph{less} than the
smallest individual resistor (1.1 k$\Omega$, say, versus the smallest resistor at 2.5
k$\Omega$), while a series combination's equivalent resistance is the full sum (10.5
k$\Omega$). Since smaller $R_{\text{eq}}$ means more power at fixed voltage, the parallel
arrangement should deliver far more power.}

\solve{
\[
R_{\text{eq,parallel}} = \left(\frac{1}{2.5\text{ k}\Omega}+\frac{1}{3.5\text{
k}\Omega}+\frac{1}{4.5\text{ k}\Omega}\right)^{-1} = 1.10\text{ k}\Omega
\;\Rightarrow\;
P_{\text{parallel}} = \frac{(100\text{ V})^2}{1101\ \Omega} = 9.1\text{ W}
\]
\[
R_{\text{eq,series}} = 2.5\text{ k}\Omega+3.5\text{ k}\Omega+4.5\text{ k}\Omega =
10.5\text{ k}\Omega
\;\Rightarrow\;
P_{\text{series}} = \frac{(100\text{ V})^2}{10{,}500\ \Omega} = 0.95\text{ W}
\]
\[
\frac{P_{\text{parallel}}}{P_{\text{series}}} = \frac{9.1\text{ W}}{0.95\text{ W}} =
\boxed{9.5}.
\]}

\review{Because both cases share the same $\Delta V$, the power ratio is exactly the
resistance ratio flipped: $R_{\text{eq,series}}/R_{\text{eq,parallel}} = 10{,}500/1101
\approx 9.5$. The parallel combination dissipates almost ten times more power from the
same supply --- a direct consequence of a parallel network's resistance always dropping
below its smallest branch, while a series network's only ever grows.}

\end{document}
